\documentclass[../../../uem_phy_oka_q.tex]{subfiles}

\begin{document}
    図\ref{fig:uem_phy_oka_2025_uni_01_q_01}のように，点Pを中心とする半径$r$の%
    半円形でなめらかな面のくぼみを持つ質量$M$の台が，なめらかで水平な床の上に置かれており，%
    台の動きがストッパーで止められている．くぼみの左端に質量$m$の小球を置いて静かに手を放すと，%
    小球は，くぼみの内面に沿って動き，くぼみの中で往復運動を始めた．%
    小球が最初に最下点Qを通過するときについて，以下の問いに答えよ．ただし，%
    重力加速度の大きさを$g$とし，速度に関しては，図の右方向を正として床に固定された$x$軸を%
    基準とせよ．また，台は左右対称の形状であり，その対称軸は$x=0$にあるとする．\\

    \begin{enumerate}
        \item%
            小球の速度を求めよ．$\{M,\,m,\,r,\,g\}$
        \item%
            小球が台のくぼみの内面から受ける垂直抗力の大きさを求めよ．$\{M,\,m,\,g\}$
    \end{enumerate}

    次に，図\ref{fig:uem_phy_oka_2025_uni_01_q_02}のようにストッパーをはずして，%
    左端に置いた小球から静かに手を放すと，小球はなめらかなくぼみの内面に沿って%
    動くとともに台も動きだし，両者は往復運動を開始した．%
    小球が最初に最下点Qを通過するときについて，以下の問いに答えよ．

    \begin{enumerate}[resume]
        \item%
            小球と台からなる物体系について水平方向の運動量保存則を適用し，%
            台の水平方向の速度$V$を小球の水平方向の速度$v$を用いて表せ．$\{M,\,m,\,v\}$
        \item%
            小球の水平方向の速度を求めよ．$\{M,\,m,\,r,\,g\}$
        \item%
            台の水平方向の速度を求めよ．$\{M,\,m,\,r,\,g\}$
        \item%
            小球が台のくぼみの内面から受ける垂直抗力の大きさを求めよ．$\{M,\,m,\,g\}$
    \end{enumerate}
    \newpage
    \noindent
    \begin{minipage}{\linewidth}
        \setlength{\abovecaptionskip}{4pt}
        \centering
            \includegraphics{fig_uem_phy_oka_2025_uni_01/fig_uem_phy_oka_2025_uni_01_q_01.pdf}
            \captionof{figure}{} \label{fig:uem_phy_oka_2025_uni_01_q_01}
            \vspace{10pt}
            \includegraphics{fig_uem_phy_oka_2025_uni_01/fig_uem_phy_oka_2025_uni_01_q_02.pdf}
            \captionof{figure}{} \label{fig:uem_phy_oka_2025_uni_01_q_02}
    \end{minipage}
\end{document}